%% doku.tex — Vorlage-Dokumentation & Referenz
%% Zeigt alle Umgebungen mit echten Beispielen.
%% Inhalt: Sortieralgorithmen als roter Faden.
%% ──────────────────────────────────────────────────────────────────────────

\section{Informations- und Hinweis-Boxen}

\subsection{infobox — Infos und Definitionen}

\begin{infobox}[Definition: Algorithmus]
  Ein \textbf{Algorithmus} ist eine endliche, eindeutige Folge von Anweisungen,
  die ein Problem in endlicher Zeit löst. Jeder Algorithmus muss
  \emph{terminieren}, \emph{korrekt} und \emph{eindeutig} sein.
\end{infobox}

\begin{infobox}
  Infoboxen können auch ohne Titel verwendet werden — z.\,B. für kurze
  Hinweise oder ergänzende Informationen im Fließtext.
\end{infobox}

\subsection{warnbox — Warnhinweise und Fallstricke}

\begin{warnbox}[Off-by-one-Fehler]
  Beim Zugriff auf Array-Indizes stets prüfen, ob der Index im Bereich
  $[0, n-1]$ liegt. Ein häufiger Fehler ist der Zugriff auf Index $n$,
  was zu einem \icode{IndexOutOfBoundsException} führt.
\end{warnbox}

\subsection{merksatz — Wichtige Sätze und Regeln}

\begin{merksatz}[Master-Theorem]
  Für Rekurrenzen der Form $T(n) = a \cdot T(n/b) + f(n)$ gilt:
  \[
    T(n) = \begin{cases}
      \Theta(n^{\log_b a})       & \text{falls } f(n) = O(n^{\log_b a - \varepsilon}) \\
      \Theta(n^{\log_b a}\log n) & \text{falls } f(n) = \Theta(n^{\log_b a}) \\
      \Theta(f(n))               & \text{falls } f(n) = \Omega(n^{\log_b a + \varepsilon})
    \end{cases}
  \]
\end{merksatz}

\begin{merksatz}
  Merksätze können auch ohne Titel stehen — etwa für kurze, prägnante
  Regeln die man unbedingt im Kopf behalten muss.
\end{merksatz}

\subsection{beispiel — Konkrete Beispiele}

\begin{beispiel}[Quicksort auf {[3, 1, 4, 1, 5, 9]}]
  \textbf{Pivot} = 3 (erstes Element). Partition ergibt:
  \begin{itemize}[nosep]
    \item Links $\leq 3$: \icode{[1, 1]}
    \item Rechts $> 3$: \icode{[4, 5, 9]}
  \end{itemize}
  Rekursiver Aufruf auf beide Teile, dann zusammensetzen:
  \icode{[1, 1, 3, 4, 5, 9]}.
\end{beispiel}

\subsection{aufgabe — Übungsaufgaben}

\begin{aufgabe}[Laufzeitanalyse]
  Bestimmen Sie die durchschnittliche Zeitkomplexität von Quicksort mithilfe
  des Master-Theorems. Welche Rekurrenz beschreibt den Algorithmus?
\end{aufgabe}

\begin{aufgabe}
  Implementieren Sie Mergesort in Python und vergleichen Sie die Laufzeit
  empirisch mit Bubblesort für $n \in \{100, 1000, 10000\}$.
\end{aufgabe}

% ─────────────────────────────────────────────────────────────────────────────
\section{Strukturelemente}

\subsection{process — Schrittfolgen}

\begin{process}
  \step{Pivot-Element aus der Liste wählen (z.\,B. erstes, letztes oder mittleres)}
  \step{Liste in zwei Teile partitionieren: Elemente $\leq$ Pivot links, $>$ Pivot rechts}
  \step{Quicksort rekursiv auf den linken Teilbereich anwenden}
  \step{Quicksort rekursiv auf den rechten Teilbereich anwenden}
  \step{Ergebnis: $\text{links} + [\text{pivot}] + \text{rechts}$}
\end{process}

\subsection{procon — Vor- und Nachteile}

\begin{procon}
  \pro{Sehr schnell in der Praxis ($O(n \log n)$ im Durchschnitt)}
  \pro{In-place — kein zusätzlicher Speicher nötig}
  \pro{Cache-freundlich durch sequentiellen Speicherzugriff}
  \con{Worst-case $O(n^2)$ bei schlechter Pivot-Wahl}
  \con{Nicht stabil — gleiche Elemente können die Reihenfolge tauschen}
  \con{Rekursionstiefe kann bei großen $n$ zum Stack-Overflow führen}
\end{procon}

\subsection{konzeptkarte — Strukturierte Übersichtskarte}

Die \texttt{konzeptkarte}-Umgebung fasst Beschreibung, Tabelle und Ablauf
in einer visuellen Karte zusammen. Mit \texttt{\textbackslash kzweispaltig}
lassen sich zwei Inhalte nebeneinander setzen.

\begin{konzeptkarte}{Mergesort}
  Ein \textbf{stabiler}, rekursiver Sortieralgorithmus nach dem
  Teile-und-Herrsche-Prinzip. Zerlegt die Liste in zwei Hälften,
  sortiert diese rekursiv und \emph{mischt} sie dann zusammen.

  \kzweispaltig{%
    \begin{stdtable}{l X}{Eigenschaft & Wert}
      Zeitkomplexität & $O(n \log n)$ immer \\
      Speicher        & $O(n)$ extra \\
      Stabil          & Ja \\
      In-place        & Nein \\
    \end{stdtable}
  }{%
    \begin{process}
      \step{Liste halbieren}
      \step{Linke Hälfte rekursiv sortieren}
      \step{Rechte Hälfte rekursiv sortieren}
      \step{Beide Hälften zusammenmischen}
    \end{process}
  }
\end{konzeptkarte}

\begin{konzeptkarte}{Heapsort}
  Nutzt einen \textbf{Max-Heap}, um das größte Element zu finden und an
  die richtige Position zu setzen. Nicht stabil, aber in-place und
  garantiert $O(n \log n)$.

  \kzweispaltig{%
    \begin{stdtable}{l X}{Eigenschaft & Wert}
      Zeitkomplexität & $O(n \log n)$ immer \\
      Speicher        & $O(1)$ extra \\
      Stabil          & Nein \\
      In-place        & Ja \\
    \end{stdtable}
  }{%
    \begin{process}
      \step{Max-Heap aufbauen (\icode{heapify})}
      \step{Wurzel (Maximum) ans Ende tauschen}
      \step{Heap-Größe um 1 reduzieren}
      \step{Schritte 2–3 wiederholen}
    \end{process}
  }
\end{konzeptkarte}

% ─────────────────────────────────────────────────────────────────────────────
\section{Code-Elemente}

\subsection{codebox — Einfacher Monospace-Block}

Für kurze Code-Fragmente ohne Syntax-Highlighting: \icode{codebox}.

\begin{codebox}
Eingabe: [5, 2, 9, 1, 7]\\
Nach Pass 1: [2, 5, 1, 7, 9]\\
Nach Pass 2: [2, 1, 5, 7, 9]\\
Nach Pass 3: [1, 2, 5, 7, 9]\\
Fertig nach 4 Durchläufen.
\end{codebox}

\subsection{syntaxbox — Code mit Syntax-Highlighting}

\subsubsection{Python}

\begin{syntaxbox}[Python]
def quicksort(arr):
    if len(arr) <= 1:
        return arr
    pivot = arr[len(arr) // 2]
    left  = [x for x in arr if x < pivot]
    mid   = [x for x in arr if x == pivot]
    right = [x for x in arr if x > pivot]
    return quicksort(left) + mid + quicksort(right)

# Beispielaufruf
data = [3, 6, 8, 10, 1, 2, 1]
print(quicksort(data))  # [1, 1, 2, 3, 6, 8, 10]
\end{syntaxbox}

\subsubsection{Java}

\begin{syntaxbox}[Java]
public class Mergesort {

    public static int[] sort(int[] arr) {
        if (arr.length <= 1) return arr;
        int mid = arr.length / 2;
        int[] left  = sort(Arrays.copyOfRange(arr, 0, mid));
        int[] right = sort(Arrays.copyOfRange(arr, mid, arr.length));
        return merge(left, right);
    }

    private static int[] merge(int[] a, int[] b) {
        int[] result = new int[a.length + b.length];
        int i = 0, j = 0, k = 0;
        while (i < a.length && j < b.length)
            result[k++] = (a[i] <= b[j]) ? a[i++] : b[j++];
        while (i < a.length) result[k++] = a[i++];
        while (j < b.length) result[k++] = b[j++];
        return result;
    }
}
\end{syntaxbox}

\subsubsection{SQL}

\begin{syntaxbox}[SQL]
-- Ergebnisse nach Score sortiert (analog zu ORDER BY = Sortieralgorithmus)
SELECT student_id, name, score
FROM   exam_results
WHERE  score >= 50
ORDER  BY score DESC, name ASC
LIMIT  10;
\end{syntaxbox}

\subsection{\textbackslash icode — Inline-Code im Fließtext}

Der Befehl \icode{\textbackslash icode\{...\}} hebt Code im Fließtext hervor:
Die Funktion \icode{mergesort(arr)} hat eine Zeitkomplexität von $O(n \log n)$.
Die Variable \icode{pivot} speichert das Pivot-Element.
Fehler vom Typ \icode{NullPointerException} treten auf, wenn \icode{arr == null}.

% ─────────────────────────────────────────────────────────────────────────────
\section{Tabellen und Layout}

\subsection{stdtable — Formatierte Tabellen}

\begin{stdtable}{l l r r}{Algorithmus & Strategie & $\varnothing$ & Worst-case}
  Quicksort  & Teile-und-Herrsche & $O(n \log n)$ & $O(n^2)$      \\
  Mergesort  & Teile-und-Herrsche & $O(n \log n)$ & $O(n \log n)$ \\
  Heapsort   & Heap               & $O(n \log n)$ & $O(n \log n)$ \\
  Bubblesort & Vergleich          & $O(n^2)$      & $O(n^2)$      \\
  Timsort    & Hybrid             & $O(n \log n)$ & $O(n \log n)$ \\
\end{stdtable}

\begin{warnbox}[Spaltendefinition]
  \icode{l}, \icode{r}, \icode{c} sind fest breite Spalten. \icode{X} füllt
  den verbleibenden Platz automatisch. \icode{p\{3cm\}} ist eine
  manuell festgelegte Breite.
\end{warnbox}

\subsection{sidetables — Zwei Inhalte nebeneinander}

\sidetables{%
  \begin{stdtable}{l r}{Struktur & Zugriff}
    Array        & $O(1)$      \\
    Linked List  & $O(n)$      \\
    Hash Map     & $O(1)$ avg. \\
    BST          & $O(\log n)$ \\
  \end{stdtable}
}{%
  \begin{stdtable}{l r}{Struktur & Einfügen}
    Array        & $O(n)$      \\
    Linked List  & $O(1)$      \\
    Hash Map     & $O(1)$ avg. \\
    BST          & $O(\log n)$ \\
  \end{stdtable}
}

% ─────────────────────────────────────────────────────────────────────────────
\section{Konfiguration der Vorlage}

\begin{infobox}[Anpassung in drei Dateien]
  \textbf{1.\ config/colormeta.tex} — Farben, Schriften, Abstände (\emph{hier alles anpassen})\\
  \textbf{2.\ config/usermeta.tex} — Name, Studiengang, Matrikelnummer\\
  \textbf{3.\ config/topicmeta.tex} — Titel, Fach, Semester, Beschreibung
\end{infobox}

\begin{merksatz}[Goldene Regel]
  Farben und Schriften werden \emph{ausschließlich} in \icode{colormeta.tex}
  definiert. Nirgendwo sonst im Dokument werden Hex-Werte oder Schriftnamen
  hart kodiert — alle Umgebungen referenzieren nur Farbnamen wie
  \icode{primary}, \icode{accent}, \icode{infoframe} usw.
\end{merksatz}

\subsection{Schriften anpassen}

\begin{codebox}
\% Standard (immer in TeX Live verfügbar):\\
\textbackslash newcommand\{\textbackslash mainfont\}\{TeX Gyre Pagella\}\\
\textbackslash newcommand\{\textbackslash sansfont\}\{TeX Gyre Heros\}\\
\textbackslash newcommand\{\textbackslash monofont\}\{TeX Gyre Cursor\}\\[4pt]
\% Premium-Alternativen (Homebrew: brew install --cask font-source-serif-4):\\
\textbackslash newcommand\{\textbackslash mainfont\}\{Source Serif 4\}\\
\textbackslash newcommand\{\textbackslash sansfont\}\{Source Sans 3\}\\
\textbackslash newcommand\{\textbackslash monofont\}\{JetBrains Mono\}
\end{codebox}

\begin{aufgabe}[Template anpassen]
  Ändern Sie die Farben \icode{primary} und \icode{accent} in
  \icode{colormeta.tex} auf eigene Werte und kompilieren Sie mit
  \icode{make build}. Beobachten Sie, wie sich alle Boxen, Tabellen
  und Überschriften automatisch anpassen.
\end{aufgabe}
